% =====================================================
% CHAPITRE 7 : DÉPLOIEMENT ET REPRODUCTIBILITÉ (CRISP-DM Phase 6)
% =====================================================
\chapter{Déploiement et reproductibilité}

% ---------------------------------------------------
\section*{Introduction}
\addcontentsline{toc}{section}{Introduction}
% ---------------------------------------------------

Un modèle performant qui ne peut pas être réexécuté par quelqu'un d'autre n'a pas grande
valeur dans un contexte professionnel. Pour ce projet, le déploiement signifie rendre
l'ensemble du travail \textbf{reproductible de bout en bout}, maintenable par un futur
contributeur et extensible si les besoins évoluent. Ce chapitre décrit l'architecture
choisie, l'organisation des sorties, le script d'orchestration, les garanties de
reproductibilité, la suite de tests unitaires mise en place, et les perspectives
d'amélioration identifiées pour une éventuelle industrialisation du système.

% ---------------------------------------------------
\section{Architecture du projet}
% ---------------------------------------------------

\subsection{Structure modulaire}

Le projet est organisé en modules fonctionnels clairement délimités, chacun ayant une
responsabilité unique (\textit{separation of concerns}) :

\begin{table}[H]
\centering
\caption{Structure des modules du projet}
\label{tab:modules}
\begin{tabular}{p{4.5cm}p{8cm}}
\toprule
\textbf{Module} & \textbf{Responsabilité} \\
\midrule
\texttt{src/preprocessing/}        & Chargement, nettoyage, normalisation, split \\
\texttt{src/feature\_engineering/} & Création des 7 variables dérivées \\
\texttt{src/models/}               & Modèles ML (LR, RF) et autoencoder PyTorch \\
\texttt{src/pipeline/}             & Orchestration entraînement et inférence \\
\texttt{src/utils/}                & Métriques, évaluation, utilitaires communs \\
\texttt{src/visualization/}        & Génération de figures et graphiques \\
\texttt{src/ollama\_integration/}  & Génération d'explications LLM via Ollama \\
\texttt{src/explainability/}       & Explicabilité SHAP et LIME \\
\texttt{tests/}                    & Tests unitaires (pytest) \\
\bottomrule
\end{tabular}
\end{table}

Cette organisation garantit que chaque module peut être modifié, testé et remplacé
indépendamment --- un avantage décisif pour la maintenance et l'évolution du système.

\subsection{Organisation des sorties}

Le pipeline produit des artefacts organisés dans le répertoire \texttt{outputs/} :

\begin{itemize}
  \item \texttt{outputs/models/} : artefacts sérialisés --- poids de l'autoencoder,
    scalers, seuils optimaux, poids des classes, scores AE
    (\texttt{ae\_scores\_val/test.npy}) et tenseurs SHAP sauvegardés
    (\texttt{shap\_values\_*.npy}) ;
  \item \texttt{outputs/reports/} : rapports JSON des phases exécutées
    (\texttt{eda\_report.json}, \texttt{prep\_report.json},
    \texttt{baseline\_report.json}, \texttt{autoencoder\_report.json}) ;
  \item \texttt{outputs/figures/} : courbes de performance, matrices de confusion,
    distributions, beeswarm SHAP, projection latente de l'AE ;
  \item \texttt{outputs/anomalies\_report.csv} : liste des transactions signalées avec
    leurs scores et niveaux de risque ;
  \item \texttt{outputs/explanations.json} : explications textuelles des anomalies
    générées par Ollama.
\end{itemize}

% ---------------------------------------------------
\section{Orchestration et reproductibilité}
% ---------------------------------------------------

\subsection{Script d'orchestration}

Le script \texttt{run\_all.py} constitue le point d'entrée unique du pipeline. Il valide
d'abord les imports des modules \texttt{src/}, puis exécute séquentiellement les six
notebooks (\texttt{01} à \texttt{06}) via \texttt{nbconvert}. Il propose plusieurs
modes d'exécution :

\begin{itemize}
  \item \textbf{Exécution complète} : lance les notebooks 01 à 06 en séquence ;
  \item \textbf{Reprise partielle} : option \texttt{--from N} pour reprendre à partir
    d'un notebook spécifique ;
  \item \textbf{Notebook unique} : option \texttt{--only N} pour n'exécuter qu'une seule
    étape ;
  \item \textbf{Vérification uniquement} : option \texttt{--check-only} pour valider
    les imports sans lancer les calculs.
\end{itemize}

\subsection{Garanties de reproductibilité}

La reproductibilité est assurée par quatre mécanismes complémentaires :

\begin{enumerate}
  \item \textbf{Graines aléatoires fixes} : \texttt{RANDOM\_STATE = 42} dans tous les
    modules qui comportent une part d'aléatoire (split, SMOTE, forêt aléatoire, dropout).
  \item \textbf{Environnement virtuel} : toutes les dépendances sont centralisées dans
    \texttt{requirements.txt} avec des contraintes de versions, à installer dans un
    environnement Python isolé.
  \item \textbf{Pipeline déterministe} : toutes les transformations (scaler, encodeurs)
    sont sérialisées via \texttt{joblib} dans \texttt{outputs/models/} pour être
    réappliquées à l'inférence sans recalcul.
  \item \textbf{Séparation stricte des données} : les ensembles de validation et test
    ne sont jamais utilisés pour l'apprentissage ni les transformations statistiques.
\end{enumerate}

% ---------------------------------------------------
\section{Tests unitaires}
% ---------------------------------------------------

\subsection{Structure de la suite de tests}

Le dépôt contient une suite de tests unitaires structurée sous \texttt{tests/} et
exécutable via \texttt{pytest -q} :

\begin{table}[H]
\centering
\caption{Suite de tests unitaires du projet}
\label{tab:tests}
\begin{tabular}{p{5.5cm}p{7cm}}
\toprule
\textbf{Fichier de test} & \textbf{Modules testés} \\
\midrule
\texttt{tests/test\_preprocessing.py}  & Pipeline de préparation des données \\
\texttt{tests/test\_models.py}         & Modèles ML et autoencoder \\
\texttt{tests/test\_visualization.py}  & Génération de figures \\
\texttt{tests/test\_utils.py}          & Utilitaires et métriques \\
\texttt{tests/test\_ollama.py}         & Intégration LLM (Ollama) \\
\bottomrule
\end{tabular}
\end{table}

Les fichiers de tests constituent un squelette structuré prêt à être alimenté. L'étape
suivante consiste à implémenter les cas de test dans ces fichiers pour atteindre
l'objectif de couverture à 80\,\%.

% ---------------------------------------------------
\section{Perspectives d'amélioration}
% ---------------------------------------------------

\subsection{Améliorations techniques}

Les résultats obtenus ouvrent plusieurs pistes d'amélioration algorithmique :

\begin{itemize}
  \item \textbf{Modèles de gradient boosting} : intégration de LightGBM ou XGBoost,
    qui constituent l'état de l'art sur les données tabulaires et pourraient améliorer
    les performances supervisées tout en restant interprétables via SHAP.
  \item \textbf{Autoencoder variationnel (VAE)} : passage d'un autoencoder déterministe
    à un VAE qui modélise la distribution des données normales sous forme probabiliste
    --- plus robuste pour détecter les anomalies de bord de distribution.
  \item \textbf{Détection en temps réel} : portage du pipeline vers une architecture
    streaming (Apache Kafka + Apache Flink) pour analyser les transactions au fil de
    l'eau, plutôt qu'en batch.
  \item \textbf{Apprentissage en ligne} : mise à jour continue des modèles au fil des
    nouvelles transactions validées par les auditeurs, pour s'adapter aux évolutions
    des comportements frauduleux.
\end{itemize}

\subsection{Améliorations métier}

Sur le plan fonctionnel, plusieurs extensions rendraient le système directement opérationnel
en production :

\begin{itemize}
  \item \textbf{Interface utilisateur} : tableau de bord (dashboard) pour les analystes
    permettant de consulter les alertes, de les valider ou de les rejeter, et de consulter
    les explications en langage naturel sans accéder au code.
  \item \textbf{Boucle de rétroaction (\textit{feedback loop})} : intégration des
    décisions humaines (fraude confirmée ou infirmée) pour ré-entraîner périodiquement
    les modèles et améliorer leur précision sur les données métier réelles.
  \item \textbf{Détection comportementale} : profiling par utilisateur ou par compte
    pour détecter non plus des transactions isolées anormales, mais des dérives
    comportementales dans le temps (ex : montants progressivement croissants, changements
    d'habitudes de transfert).
\end{itemize}

% ---------------------------------------------------
\section*{Conclusion}
\addcontentsline{toc}{section}{Conclusion}
% ---------------------------------------------------

Ce chapitre a présenté l'architecture de déploiement et les mécanismes de reproductibilité
du projet. La \textbf{structure modulaire} et l'orchestration via \texttt{run\_all.py}
fournissent une base maintenable et facilement reprise par un nouveau développeur. Les
quatre garanties de reproductibilité --- graines fixes, environnement versionné, artefacts
sérialisés, séparation étanche des données --- assurent que les résultats présentés dans
ce rapport sont intégralement reproductibles. La suite de tests structurée pose les bases
d'un contrôle qualité progressif. Les perspectives identifiées offrent de nombreuses pistes
pour prolonger ce travail vers un \textbf{système de production opérationnel}, capable
d'analyser des flux transactionnels en temps réel et d'apprendre de façon continue
des retours des auditeurs.
